\appendix

\section{Appendix}
\label{sec:appendix}

\subsection{Result sources}

The manuscript uses the following summary files as the source of truth. The main benchmark summary contains 120 model-level rows: six datasets, four operators, five model families, and five-fold aggregation for each combination.
\begin{itemize}
    \item \texttt{records/LATEST/summaries/benchmark\_summary.csv}
    \item \texttt{records/LATEST/summaries/branch\_ablation\_summary.csv}
    \item \texttt{records/LATEST/summaries/parameter\_sensitivity\_summary.csv}
    \item \texttt{records/LATEST/summaries/tuned\_candidate\_summary.csv}
    \item \texttt{records/LATEST/summaries/mechanism\_compact\_summary.csv}
    \item \texttt{records/SYN\_MULTI\_FULL/summaries/benchmark\_summary.csv}
\end{itemize}

\subsection{Supplementary diagnostics}

Additional fold-level rows, CKA summaries, residual traffic summaries, and full sensitivity tables are retained in the repository records directory for supplementary conversion.

\subsection{Evidence map}

Table~\ref{tab:evidence_map} summarizes where each main empirical claim is supported in the manuscript. This compact table is included to keep the main text readable while still making the result-to-claim link explicit.

\begin{table}[h]
\centering
\small
\caption{Compact map from empirical claims to manuscript evidence.}
\label{tab:evidence_map}
\begin{tabular}{p{0.42\linewidth} p{0.44\linewidth}}
\toprule
Claim & Primary evidence \\
\midrule
MatrixRes is the most frequent winner in the full benchmark but not a universal winner. & Table~\ref{tab:operator_wins}, Table~\ref{tab:winners_by_operator}, and Figure~\ref{fig:model_win_counts}. \\
GCNConv results remain dataset-dependent under a fixed operator. & Table~\ref{tab:main_benchmark} and Figure~\ref{fig:main_benchmark}. \\
Branch count produces a rise-then-fall pattern rather than monotonic improvement. & Table~\ref{tab:branch_best} and Figure~\ref{fig:branch_ablation}. \\
Useful branch expansion is associated with separation without excessive residual traffic. & Table~\ref{tab:mechanism_snapshot} and Figure~\ref{fig:mechanism}. \\
Fold-0 sensitivity should not be treated as final evidence. & Table~\ref{tab:sensitivity}, Figure~\ref{fig:sensitivity}, and Table~\ref{tab:tuned_candidates}. \\
Matrix-style reuse becomes more favorable in higher-class synthetic structural discrimination. & Table~\ref{tab:synthetic_class_scaling}, Figure~\ref{fig:synthetic_multiclass_scaling}, Table~\ref{tab:synthetic_high_class}, and Table~\ref{tab:synthetic_distribution_high_class}. \\
\bottomrule
\end{tabular}
\end{table}

\subsection{Synthetic distribution-level summary}

Table~\ref{tab:synthetic_distribution_high_class} reports the higher-class synthetic results by graph family. The table is placed in the appendix because the main text focuses on the class-count scaling trend. It shows that MatrixRes is strongest on BA, ER, and random-regular graphs, while SBM and WS remain less favorable to matrix-style reuse.

\begin{table}[h]
\centering
\scriptsize
\caption{Synthetic high-class accuracy by graph family, averaged over \(C=4,\ldots,8\) and four operators.}
\label{tab:synthetic_distribution_high_class}
\begin{tabular}{l c c c c c}
\toprule
Family & Plain & VerticalRes & HorizontalRes & MatrixRes & MatrixResGated \\
\midrule
BA & 0.7874 & 0.8177 & 0.7855 & \(\mathbf{0.8271}\) & 0.8257 \\
ER & 0.8777 & 0.8756 & 0.8728 & \(\mathbf{0.8807}\) & 0.8796 \\
Regular & 0.7505 & 0.7923 & 0.7739 & \(\mathbf{0.8077}\) & 0.8057 \\
SBM & 0.6268 & \(\mathbf{0.6286}\) & 0.6218 & 0.6264 & 0.6176 \\
WS & \(\mathbf{0.3709}\) & 0.3511 & 0.3693 & 0.3202 & 0.3430 \\
\bottomrule
\end{tabular}
\end{table}
